\subsection{Integration systems satisfying the monotone convergence theorem}

\begin{definition}\label{def:satisfies-mct}
  We say that an integration system $(X,L,I)$ on a set $X$
  \textbf{satisfies the monotone convergence theorem} if, for every function
  $f:X\to\mathbb{R}$ and every monotonically nondecreasing sequence
  $\{f_n\}_{n=1}^{\infty}\subset L$,
  \[
    f_n(x)\uparrow f(x)\quad(\forall x\in X),\qquad\sup_n I(f_n)<+\infty
  \]
  implies
  \[
    f\in L,\qquad I(f)=\lim_{n\to\infty} I(f_n) .
  \]
\end{definition}

\noindent
This condition is a restatement in axiomatic form of the classical theorem of
Beppo Levi (1906), and \cite{CWS} calls a triple satisfying it \textbf{closed}
(its Definition 1.4.1).
The correspondence of terms with that book is collected in
Remark~\ref{rem:extended-real-collapse}.

\begin{definition}
  The full subcategory of the category $\mathbf{Int}$ of integration systems
  consisting of the integration systems that satisfy the monotone convergence
  theorem is written
  \[
    \mathbf{MInt}
  \]
  and called the \textbf{category of integration systems satisfying the monotone
  convergence theorem}.
\end{definition}
